\chapter{A research program beyond bounded mass}
\label{chap:research-program}

The preceding chapters establish several parts of the theory on every compact Hermitian manifold. They also show why a complete generalization of \cite{SGPT} cannot be obtained by changing the background notation in its statements. We now collect the remaining problems in a form that separates established literature, missing proofs in this manuscript, and proposed mathematical statements.

A question below records a target for this program, without claiming that it has never appeared elsewhere. A working conjecture is a proposed answer supported by the indicated cases; it is not used in any proof in the book. The absence of a bounded mass assumption remains part of each question unless an additional regime is stated explicitly.

\section{Three different kinds of finiteness}

Let $X$ be a connected compact complex manifold, let $\theta$ be a smooth real $(1,1)$-form on $X$, and let $\varphi\in\PSH(X,\theta)$ be a reference potential. Write $\mathcal E_{\mathrm{env}}=\mathcal E_{\mathrm{env}}(X,\theta,\varphi)$ for the envelope-defined class of \cref{def:relative-envelope-class}.

The examples force a distinction among finite linear trace, finite non-pluripolar top-degree mass, and finite equation normalization. Every potential has a finite nonnegative total trace against a fixed Gauduchon metric; for a Hermitian background this total is strictly positive. Its extended Monge--Amp\`ere measure can nevertheless have infinite mass. The smooth equation still has a finite positive normalization, and the Gauduchon formula extends that number to a finite nonnegative quantity on all pseudoeffective smooth backgrounds.

These statements do not describe successive versions of one invariant. The trace determines the whole potential up to a constant by \cref{thm:gaud-trace-homeomorphism}. The non-pluripolar measure discards mass on pluripolar sets. The equation normalization is invariant under smooth changes of representative of the background. Each retains different information. The construction in \cref{prop:integrable-infinite-mass,cor:strong-linear-infinite-mass} makes the distinction stronger: the full currents have integrable coefficients and their traces converge in total variation, while even the pointwise determinant densities have infinite integral. Linear control alone therefore cannot supply the nonlinear integrability needed for an equation.

\Needspace{8\baselineskip}
The distinction also affects the meaning of a full mass class. Without bounded mass, the envelope-defined class $\mathcal E_{\mathrm{env}}$ can contain a potential of infinite extended mass, as \cref{prop:goodness-infinite-mass} shows. Thus membership in this class cannot replace a finiteness hypothesis when a proof integrates its Monge--Amp\`ere measure. In relative zero-domination, finiteness is proved only for the auxiliary measure to which finite-measure domination is applied. Other positive integrals in the capacity argument are allowed to be infinite.

\section{General measures and the remaining hypotheses}

Let $X$ be a connected compact complex manifold of dimension $n$, let $h$ be a Hermitian form, and fix a smooth positive probability volume form $\mathrm dV$. Smooth real $(1,1)$-backgrounds are denoted by $\theta$; their additional hypotheses are specified in each regime below. For a semipositive background, $v_-$ and $v_+$ denote the infimum and supremum of the masses of bounded potentials. For nonzero $0\leq f\in L^1(\mathrm dV)$, let $b_h(f)$ be the greatest coefficient of a genuine subsolution $\mathrm{MA}^{\mathrm{ext}}_h(u)\geq b_h(f)f\,\mathrm dV$, as in \cref{thm:individual-density-subsolution}. No positivity of this coefficient is assumed.

Two general-measure regimes have been reconstructed. For a semipositive big background, \cref{sec:general-measure-relative} proves \cref{thm:als-general-measure} of \cite{ALS26}, with its model condition and $v_-(\theta)>0$, without bounded mass. It covers the positive-exponent equation and the zero-exponent equation with its unique positive scalar. For an arbitrary smooth big background, \cref{thm:big-model-exponential,thm:big-model-normalized} proves both equations under bounded mass and positive lower volume of the model reference. The latter hypothesis implies metric positive lower volume in this regime, as explained after \cref{thm:li-general-measure}. \Cref{cor:big-model-lp-type} gives the prescribed ordinary type for $L^p$ densities, whereas the minimal ordinary-type $L^p$ theorem in \cref{sec:big-minimal-existence} needs neither metric volume property.

The two proofs use different compactness estimates. In the semipositive argument, relative zero-domination constructs an auxiliary envelope whose measure is supported where the envelope is nonnegative. Locality with one bounded truncation proves that this auxiliary measure is finite; the limiting lower roof need not have finite mass. Contact support, weighted lower semicontinuity under capacity convergence and a common exponential envelope give the moving-measure limit. Uniform sublevel tails then prove finiteness of the solution measure. The arbitrary-background argument uses an explicit rooftop between different classes and finite upper relative volumes for compactness and weighted-mass maximization. Removing bounded mass requires replacing these estimates and still proving finiteness of the solution measure; existence of the formal envelopes alone is insufficient.

Universal solvability cannot omit the lower-volume condition. In a Hermitian background, \cref{cor:positive-lower-mass-solvability} characterizes solvability for every finite non-pluripolar measure by $v_-(h)>0$. When it vanishes, \cref{thm:lower-mass-obstructing-measure} gives an $L^1$ probability density, bounded below by a positive constant, which admits no positive scaled subsolution among genuine potentials. These obstructing densities are $L^1$ dense, although every smooth strictly positive density is solvable. \Cref{thm:vanishing-lower-mass} realizes this obstruction on a compact threefold.

For an individual nonnegative $L^1$ density, the relevant number is the attained subsolution threshold $b_h(f)$ of \cref{thm:individual-density-subsolution}. It is the decreasing limit of the bounded-truncation normalization constants and is positive exactly when some positive multiple of $f\,\mathrm{d}V$ is dominated by an extended Monge--Amp\`ere measure. The upper-envelope and damped-continuity proof gives a subsolution, without asserting an exact equation or finite mass. Smooth tests alone do not suffice, by \cref{thm:smooth-tests-insufficient}. The Gauduchon formula in \cref{thm:integrable-gauduchon-formula} makes the threshold upper semicontinuous in $L^1$ and identifies the Banach space of admissible density roots. This space is independent of the Hermitian metric by \cref{prop:individual-density-banach,prop:individual-threshold-backgrounds}. Positive lower mass makes its domain all of $L^n$; finite upper mass makes the root norm complete, by \cref{prop:individual-root-completeness}. These are distinct properties.

\begin{problem}\label{qu:individual-equation-threshold}
If $0\leq f\in L^1(\mathrm{d}V)$ is nonzero and $b_h(f)>0$, does there exist $u\in\PSH(X,h)$ such that
\[
 \mathrm{MA}^{\mathrm{ext}}_h(u)=b_h(f)f\,\mathrm{d}V?
\]
Determine the singularity or envelope class in which such an equation should be posed.
\end{problem}

There are two established general cases. Under $v_-(h)>0$, \cref{cor:positive-lower-mass-exact-threshold} gives an exact solution in $\mathcal E_{\mathrm{env}}(X,h,0)$ for every nonzero nonnegative $L^1$ density. Finite-measure comparison identifies its scalar, and dominated continuity gives capacity convergence along every $L^1$-convergent subsequence of normalized greatest exponential solutions. This includes the Hopf example without bounded mass. For $L^p$ densities, $p>1$, the continuous equation and \cref{prop:individual-normalizer-lp} give the answer without positive lower mass.

The remaining case in \cref{qu:individual-equation-threshold} has $v_-(h)=0$ and $b_h(f)>0$. \Cref{cor:individual-positive-threshold-exponential} gives greatest exponential solutions, convergence of their scalars to $b_h(f)$ and total-variation convergence of their measures. Capacity convergence of the normalized potentials is the missing step. Under the root-tail condition of \cref{prop:root-uniform-tails}, \cref{prop:exponential-root-tail-modulus} quantifies the scalar and measure errors; a rate linear in the exponent needs additional control. Positivity and integrability of $f$ alone do not ensure a positive threshold. Exponential subsolutions when $b_h(f)=0$, and an extension of the density criterion to arbitrary non-pluripolar measures, are separate questions.

Gauduchon dual attainment gives several sufficient regimes. Write $\mathcal D_h$ for the finite-root-norm space of \cref{prop:individual-density-banach} and $\mathcal D_h^0$ for the closure of bounded functions in that norm. If the density root $f^{1/n}$ lies in $\mathcal D_h^0$, \cref{thm:integrable-current-attainment,prop:integrable-threshold-adjoints} give a maximizing current and represent all supporting probabilities. These probabilities are absolutely continuous, and the trace errors along smooth maximizing sequences tend to zero in total variation by \cref{prop:threshold-trace-saturation}. Uniform integrability along just one maximizing sequence already gives a maximizing current and convergence of its determinant-weighted measures, by \cref{prop:selected-gauduchon-attainment}. This weaker hypothesis on the chosen sequence does not establish the representation of all supporting probabilities.

In dimension at least two, a smooth strictly positive maximizer forces every threshold subsolution to be smooth and exact, with uniqueness modulo constants, and makes the density root smooth and nonnegative; see \cref{cor:smooth-dual-exact-equation}. For a closed background, \cref{cor:kahler-smooth-dual-attainment} characterizes smooth dual attainment by smooth strict positivity of the density. The strict positivity conclusion fails for nonclosed backgrounds on every compact complex manifold of dimension at least two, by \cref{ex:nonclosed-dual-zero-density}.

\Cref{thm:semipositive-dual-exact-equation} also allows a smooth semipositive maximizer which is strictly positive outside a proper analytic set $Z$. It gives an exact global equation, smoothness off $Z$ and uniqueness of all threshold subsolutions modulo constants. The proof takes the maximum of two subsolutions and uses a Gauduchon integral to eliminate their divisorial difference; no singularity type is prescribed. The positive pole-divisor classes are linearly independent by \cref{cor:semipositive-dual-divisor-independence}. Write $b=b_h(f)$ and, for the exponential parameter $\lambda>0$, let $u_\lambda$ be the greatest exponential solution and put $c_\lambda=\exp(\lambda\sup_Xu_\lambda)$. In this regime \cref{prop:semipositive-dual-exponential-limit} gives whole-family normalized $L^1$ convergence, a fixed lower barrier and, for a fixed finite constant $C$, the estimate $b\leq c_\lambda\leq b\mathrm{e}^{C\lambda}$. A genuine smaller-background envelope of the threshold solution gives capacity convergence by \cref{cor:semipositive-dual-capacity-limit}; $L^p$ data with $p>1$ give a continuous solution and uniform convergence by \cref{cor:semipositive-dual-lp-limit}. \Cref{ex:semipositive-projective-maximizer} has a smooth solution but a nondifferentiable density root on $Z$. For general $L^1$ data, automatic genuineness of that envelope and regularity across $Z$ remain open here.

The moment criterion \cref{thm:gauduchon-moment-threshold} uses no regularity of a maximizing current. Uniform first moments of normalized quasi-plurisubharmonic functions against one smooth maximizing sequence give a fixed lower barrier for the exponential solutions. Their normalized cluster limits lie below the greatest normalized threshold subsolution $\Phi_f$. If $P_{\delta h}(\Phi_f)$ is genuine for some $0<\delta<1$, the lower-barrier criterion gives subsequential capacity convergence and an exact threshold equation, without either mass property. An entropy bound on the test probabilities supplies both the first moments and the selected-sequence uniform integrability above. By \cref{prop:hermitian-dual-entropy-bounds}, such an entropy bound is equivalent to a uniform bound on these probabilities for $H_0$, the convex dual of the zero-reference normalization defined in \cref{chap:normalization}. It supplies no smaller-background envelope.

Thus singular-current attainment alone does not settle the equation. At zero lower mass, it is still unknown here whether root tails suffice, or even whether $\mathcal D_h^0=\mathcal D_h$. \Cref{prop:root-tail-continuity} characterizes the tail condition by continuity under uniformly bounded logarithmic perturbations tending to zero in measure. This is an order-interval statement, not unrestricted $L^1$ continuity of the threshold.

Two cautions about the existence proofs remain essential. The prescribed-type carrier of \cref{prop:relative-measure-carrier} can be unbounded and have infinite extended mass, so it cannot directly invoke the bounded subsolution theorem. The semipositive proof instead obtains its comparison potential from the minimal-reference equation and compares it using reference zero. Nor can an $L^1$ limit be inserted directly into a Monge--Amp\`ere equation: the finite atomic curve models of \cref{thm:curve-ideal-lone-completion} converge to a smooth positive-mass potential while their non-pluripolar measures vanish, and the measure of a maximum can acquire interface mass. The existence proofs therefore derive capacity convergence from their specific measure control and use the proved contact inequalities on the appropriate pieces of lower envelopes.

\Needspace{15\baselineskip}
Interior regularity in an arbitrary smooth big background is treated in \cref{thm:nonclosed-interior-smoothness}: a locally H\"older solution with minimal singularities is smooth where a strict analytic-singularity barrier is smooth and the density is smooth and positive. The global $L^p$ case follows from \cref{thm:big-minimal-holder,cor:big-lp-interior-smoothness}, whose proofs use neither metric mass property.

More generally, \cref{cor:big-reference-uniform-convergence} gives compactness for uniform convergence of differences in a fixed ordinary type above the strict barrier, under common reference bounds and a uniform $L^p$ density bound. This compactness makes no regularity assertion for a nonminimal type. For the minimal type, \cref{prop:big-exponential-clusters} identifies every normalized exponential cluster equation and proves convergence of the scalar and the densities. Uniqueness within a symmetry-invariant class gives whole-family potential convergence by \cref{cor:big-exponential-symmetries}; general selection remains separate from this sufficient condition.

For a compact K\"ahler manifold and a smooth positive global volume form,
\cref{thm:big-class-ample-smoothness} proves smoothness of the full-mass
solution on the ample locus of every big class, including non-nef classes.
Smoothness of the density only on this open set does not imply global $L^p$ integrability. On an arbitrary smooth-density open subset, it does not even imply continuity of a bounded minimal solution: \cref{thm:compact-local-density-discontinuity} gives a compact K\"ahler counterexample with constant density on a ball. Under local closedness, \cref{thm:smooth-boundary-trace-regularity} gives a sufficient condition in terms of smooth pointwise boundary traces. Establishing the H\"older premise from density information on the entire regular locus of the strict barrier, treating a more general prescribed singularity type, and understanding the behavior across the barrier's pole set remain separate questions.

\section{Positive Gauduchon volume and bigness}

Let $X$ be a connected compact complex manifold of dimension $n$, fix a Hermitian form $h$ and a smooth strictly positive probability form $\mu$, and let $\theta$ be a smooth real $(1,1)$-background. Use $\lambda_\mu(\theta)$ and $\mathcal V_\mu(\theta)=\lambda_\mu(\theta)^n$ from \cref{def:gaud-variational-volume}. Nefness means that $\theta$ admits smooth representatives bounded below by $-\epsilon h$ for every $\epsilon>0$. Write $\lambda_*(\theta)$ for the infimum of $\lambda_\nu(\theta)$ over smooth strictly positive probability forms $\nu$.

\begin{conjecture}\label{conj:nef-gaud-volume-big}
Let $\theta$ be a nef smooth real $(1,1)$-background on a compact complex manifold. Then
\[
 \mathcal V_\mu(\theta)>0
 \quad\Longleftrightarrow\quad
 \theta\text{ is big}.
\]
Here $\mu$ is any smooth strictly positive probability measure, and bigness means the existence of a genuine potential $u$ and $a>0$ with $\theta+\ddc u\geq ah$ for a Hermitian metric $h$.
\end{conjecture}

The choice of $\mu$ does not affect the proposed equivalence, by \cref{prop:gaud-volume-density-comparison}. It is essential to keep one such density fixed. Replacing $\lambda_\mu$ by its infimum $\lambda_*$ over all smooth probability densities would make the implication from bigness false: the Hermitian backgrounds of \cref{thm:vanishing-lower-mass} are big but have $\lambda_*=0$. Bigness implies positivity: if $\theta-ah$ is pseudoeffective, monotonicity gives $\lambda_\mu(\theta)\geq a\lambda_\mu(h)>0$. In dimension one the converse follows from the integral formula. For closed nef classes on a compact K\"ahler manifold, the volume is the usual top self-intersection, so the statement agrees with the standard nef-and-big criterion recalled in \cite{SGPT}. The issue here is the converse for nonclosed backgrounds on arbitrary compact complex manifolds.

\Needspace{8\baselineskip}
\subsection{Collapse of normalized Gauduchon tests}
\label{sec:gauduchon-collapse}

The root volume and the bigness threshold divide the same pairing by different quantities. In the notation $A_\theta$, $A_h$ of \cref{sec:gaud-volume-bigness-threshold}, \cref{prop:gaud-positive-volume-collapse} shows that, for $n\geq2$, a background of positive root volume that is not big forces $D_\mu(\gamma_j)\to0$ along every sequence of Gauduchon metrics with $A_h(\gamma_j)=1$ and $A_\theta(\gamma_j)\to0$, and that a positive lower bound for $D_\mu$ along one sequence minimizing the threshold would give bigness. For the radial background of \cref{thm:radial-positive-volume-nonbig}, the normalized $(n-1)$st powers of a family of balanced test metrics (\cref{def:balanced-pluriclosed}) converge weakly to the integration current of the degenerate fiber (\cref{prop:radial-localized-balanced-tests}), and the volume-minimizing family approaches the bigness threshold only in dimension three, by \cref{prop:radial-volume-test-limits}.

The quotient $D_\mu/A_h$ is always bounded above, by \eqref{eq:gaud-collapse-basic} with $\theta=h$. What is missing for the converse direction is a lower bound for $D_\mu$ along suitable sequences minimizing the threshold. Compactness of normalized potentials does not give compactness of these Gauduchon metrics.


The collapse criterion suggests a possible use of the Gauduchon equation: replace a degenerating sequence of tests by metrics with controlled volume density while retaining the relevant pairings. There is a serious constraint. If $\theta$ is closed, pairing against $\theta$ only depends on the Aeppli class of the $(n-1,n-1)$-test form. For nonclosed $\theta$, integration by parts leaves derivatives of $\theta$, so an Aeppli-preserving deformation need not retain this pairing. The smooth Gauduchon equation is available; the required preservation or control of the numerator is a separate estimate.

This is why a proof should track the full smooth background and its derivatives. A cohomological argument that silently discards these terms would establish only the closed special case.

\section{Connections with non-K\"ahler geometry}
\label{sec:non-kahler-geometric-problems}

The general theory should explain geometric obstructions as well as solve equations. We describe six connections. In each case the distinction between a closed Bott--Chern class and an arbitrary smooth background is essential. The literature statements below are recorded with their hypotheses; the subsequent questions ask what the constructions of this book could add.

\subsection{Nef classes, bigness and the Fujiki class}

Let $X$ be a connected compact complex manifold of dimension $n$, fix a Hermitian form $h$ and a smooth strictly positive probability form $\mu$, and let $\alpha$ denote a closed real $(1,1)$-form or its Bott--Chern class. Integrals of closed classes are computed using smooth representatives. Use $v_+(h)$ for the bounded-potential mass supremum, $\lambda_\mu$ and $\mathcal V_\mu$ from \cref{def:gaud-variational-volume}, and $D_\mu$ from \cref{sec:gauduchon-collapse}.

A compact complex manifold belongs to the Fujiki class $\mathcal C$ if it is bimeromorphic to a compact K\"ahler manifold. The classical characterization by the existence of a closed K\"ahler current is due to Demailly--P\u aun. Their transcendental Grauert--Riemenschneider conjecture asks whether a closed nef real $(1,1)$-class $\alpha$ with
\[
 \int_X\alpha^n>0
\]
is big on an arbitrary compact complex manifold. Equivalently in this situation, the proposed conclusion is $X\in\mathcal C$; the source statements are recorded in \cref{sec:history-geometry}.

The bounded mass property is closely tied to this problem. Under exactly the preceding nefness and positive top-intersection hypotheses, Guedj--Lu prove
\[
 \alpha\text{ is big}\quad\Longleftrightarrow\quad v_+(h)<\infty
\]
for a Hermitian metric $h$. Boucksom--Guedj--Lu also characterize $\mathcal C$ by bounded mass together with positivity of their upper class volume for some closed Bott--Chern class. That upper volume is not the invariant $\mathcal V_\mu$ defined here.

The failure of bounded mass on the Hopf threefold does not refute the Demailly--P\u aun conjecture. A counterexample would also need a closed nef class with positive top self-intersection. Conversely, a proof of the conjecture cannot impose bounded mass as an unexplained general fact. The Hermitian framework must account for why that property follows in the particular geometry forced by such a class.

The elliptic threefold of \cref{thm:vanishing-lower-mass} supplies the complementary test: it has bounded mass, but no big closed class. Its entire nef cone is the nonnegative span of two smooth rank-one classes, and every top self-intersection is zero, by \cref{prop:elliptic-bott-chern-cones}. Hence this lower-mass example also leaves the geometric conjecture untouched. The exact boundary calculation in \cref{prop:elliptic-boundary-volume} describes its degeneration without identifying a nonclosed metric with a closed class.

The volume of this book separates two possible steps:
\begin{equation}\label{eq:dp-two-bridges}
 \int_X\alpha^n>0
 \quad\stackrel{?}{\Longrightarrow}\quad
 \mathcal V_\mu(\alpha)>0
 \quad\stackrel{?}{\Longrightarrow}\quad
 \alpha\text{ is big}.
\end{equation}
Neither arrow is established here in the unrestricted compact complex setting. The second is the closed special case of \cref{conj:nef-gaud-volume-big}. On a compact K\"ahler manifold both follow from the usual nef volume formula and the classical nef-and-big criterion. Without closedness, positivity of $\mathcal V_\mu$ cannot characterize $\mathcal C$: every Hermitian metric $h$ already satisfies $\mathcal V_\mu(h)>0$, on every compact complex manifold.

\begin{problem}\label{qu:closed-nef-gaud-bridge}
Let $\alpha$ be closed and nef with $\int_X\alpha^n>0$. Prove a positive lower bound for $\lambda_\mu(\alpha)$, and rule out the collapse in \cref{prop:gaud-positive-volume-collapse} for a sequence minimizing the bigness threshold of $\alpha$.
\end{problem}

This is a proposed route to the geometric conjecture, not a claimed weaker substitute already known to suffice automatically. The relevant test metrics can be normalized by $\int_Xh\wedge\gamma^{n-1}=1$. A positive lower bound for $D_\mu(\gamma)$ along one threshold-minimizing sequence would finish the second step. The Gauduchon equation can change the volume density while retaining an Aeppli class. Since $\alpha$ is closed, it also retains the numerator paired against $\alpha$. The missing estimate is control of the trace normalization and determinant denominator after that change. Merely prescribing a volume density does not give this estimate.

\subsection{Positive lower volume and concentration}

Let $X$ be a connected compact complex manifold of dimension $n$ and let $h$ be a Hermitian form. Its bounded-potential mass infimum and supremum are $v_-(h)$ and $v_+(h)$. For $n\geq2$, write $\mathscr G_h=\{T\geq0:\ T\text{ has bidegree }(n-1,n-1),\ \ddc T=0,\ \int_Xh\wedge T=1\}$. The measure $\Gamma_0(T)$ is the determinant-weighted measure of \eqref{eq:weak-gaud-functional} with logarithmic weight zero and the volume form chosen below.

Positive lower mass fails in general. The threefold in \cref{thm:vanishing-lower-mass} has $v_-(h)=0$ for every Hermitian metric, while \cref{prop:lower-mass-balanced-metric} supplies a balanced metric with finite upper mass. The Hopf threefold gives the opposite combination, $v_-(h)=6$ and $v_+(h)=\infty$ in its stated normalization. The upper and lower mass properties are independent.

For $n\geq2$, the determinant measures give an exact analytic formulation. Fix a smooth positive probability measure $\mathrm{d}V$, and write
\[
 \Gamma_0(T)=k_T\,\mathrm{d}V,\qquad T\in\mathscr G_h.
\]
By \cref{thm:gaud-determinant-endpoint},
\begin{equation}\label{eq:pvp-endpoint-geometric}
 v_-(h)>0
 \quad\Longleftrightarrow\quad
 \sup_{T\in\mathscr G_h}
       \|k_T\|_{L^{n/(n-1)}(\mathrm{d}V)}<\infty.
\end{equation}
The exact supremum is $v_-(h)^{-1/n}$, including the value $+\infty$. Uniform estimates for every smaller exponent $1<q<n/(n-1)$ hold without either volume property. The lower-mass example violates the endpoint already for smooth strictly positive tests. Thus positivity, $\ddc T=0$ and normalized trace alone do not control concentration at that exponent.

\begin{problem}\label{qu:geometric-gaud-endpoint}
Find geometric conditions that imply the endpoint estimate in \eqref{eq:pvp-endpoint-geometric}, and identify how they exclude the concentration mechanism of \cref{sec:vanishing-lower-mass}. For a fixed smooth density, determine whether the smaller family of maximizing Gauduchon currents admits stronger estimates than the full normalized cone.
\end{problem}

The balanced condition and bounded upper mass are insufficient, even together. The explicit identity \eqref{eq:lower-mass-vertical-energy} relates the failing lower bound to a vertical Dirichlet integral on an elliptic fibration. Meanwhile the horizontal curvature forms $\alpha,\beta$ of \cref{sec:vanishing-lower-mass} have a positive exact product $\alpha\wedge\beta$, which excludes pluriclosed metrics and closed K\"ahler currents. These calculations supply geometric information about this example without giving a classification of all lower-mass degeneration.

The equation consequence is exact for a Hermitian background: \cref{cor:positive-lower-mass-solvability} characterizes positive lower mass by solvability for all finite non-pluripolar measures. When it vanishes, even a positive $L^1$ density bounded away from zero can admit no positive scaled subsolution. The remaining equation questions concern specified measures or additional geometric hypotheses, and separately the effect of the exponential weight.

\subsection{Primitive elliptic bundles}
\label{sec:primitive-elliptic-bundles}

On a primitive elliptic bundle, the variation of the top-degree mass is a vertical Dirichlet integral. This gives the exact upper mass and a positive lower bound for separated potentials. The lower mass over all potentials remains a geometric question, since the collapse construction over a product of curves does not transfer directly to a base without curves.

Let $(B,\eta)$ be a connected compact K\"ahler surface and let $L\to B$ be a holomorphic line bundle such that
\[
 \int_Bc_1(L)\wedge\eta=0.
\]
Choose a smooth Hermitian metric on $L$ whose curvature is primitive with respect to $\eta$. Such a metric exists: for an arbitrary metric, the trace of the curvature with respect to $\eta$ has mean zero by the condition above, and the scalar Poisson equation gives a smooth change of metric that cancels this trace. Let $c$ be the resulting smooth closed curvature representative of $-c_1(L)$, so that $\eta\wedge c=0$. Write $\pi\colon L\to B$ for the bundle projection and $w$ for a nonzero vector in a fiber, with norm taken in the chosen metric. On the complement $L^\times$ of the zero section, put
\[
 t=\log|w|^2,\qquad \ddc t=\pi^*c.
\]
For $q>1$, form the quotient
\[
 X=L^\times/\langle w\mapsto qw\rangle,
 \qquad \pi\colon X\longrightarrow B,
 \qquad T=2\log q,
 \qquad A=\int_B\eta^2,
\]
where we use the same symbol $\pi$ for the induced projection. The action increases $t$ by $T$. It is free and properly discontinuous, and the compact annulus $1\leq|w|\leq q$ meets every orbit. Hence $X$ is a connected compact complex threefold; in this subsection it replaces the general manifold of the chapter. Over a trivializing open set $U\subset B$, a local holomorphic frame of $L$ identifies the fibration with $U\times(\C^*/q^{\mathbb Z})$.
\Needspace{5\baselineskip}
Since $\mathrm{d}t$ and $\dc t$ are invariant under the action, they descend to $X$, and we set
\[
 D=\mathrm{d}t\wedge\dc t,\qquad \omega=\pi^*\eta+D.
\]
The form $D$ has rank one; it is semipositive and positive on the vertical line, and its kernel projects isomorphically onto the tangent space of $B$. Hence $\omega$ is Hermitian.

\Needspace{17\baselineskip}
\begin{proposition}\label{prop:primitive-elliptic-metric}\index{Elliptic fibration!primitive bundles}
The metric $\omega$ is balanced and satisfies
\[
 \ddc\omega=-\pi^*(c^2)\geq0,\qquad
 \int_X\omega^3=3TA.
\]
For every smooth real $u$,
\begin{equation}\label{eq:primitive-elliptic-mass}
 \int_X(\omega+\ddc u)^3
 =3TA-3\int_X\mathrm{d}u\wedge\dc u\wedge\bigl(-\pi^*(c^2)\bigr).
\end{equation}
In particular $v_+(\omega)=3TA$. The canonical bundles satisfy $K_X\simeq\pi^*K_B$. If $c_1(L)\ne0$ in real cohomology, $X$ admits neither a pluriclosed metric nor a closed K\"ahler current. If $c_1(L)=0$ in real cohomology, the chosen $\omega$ is K\"ahler.
\end{proposition}
\begin{proof}
Write $a=\mathrm{d}t$ and $b=\dc t$. Then $\mathrm{d}a=0$, $\mathrm{d}b=\pi^*c$, $\dc a=-\pi^*c$ and $\dc b=0$. Consequently
\[
 D^2=0,\qquad \mathrm{d}D=-a\wedge\pi^*c,
 \qquad \ddc D=-\pi^*(c^2),
\]
\[
 \mathrm{d}(\omega^2)=-2\pi^*(\eta\wedge c)\wedge a=0,
\]
so $\omega$ is balanced. Since $c$ is primitive, its eigenvalues relative to $\eta$ are $\lambda$ and $-\lambda$ for a function $\lambda\geq0$, so $c^2=-\lambda^2\eta^2$. This proves the asserted sign of $\ddc\omega$.

To compute the volume, fix an elliptic fiber $E=\pi^{-1}(b_0)$, where $b_0\in B$. In a logarithmic fiber coordinate $x+\mathrm{i}y$, the real periods are $\log q$ and $2\pi$, while $t$ equals $2x$ up to an additive constant and $D|_E=\pi^{-1}\mathrm{d}x\wedge\mathrm{d}y$. Thus $\int_ED=T$. Since $\omega^3=3\pi^*(\eta^2)\wedge D$, the volume formula follows.

For the mass identity, expand the cube in \eqref{eq:primitive-elliptic-mass}. The terms with one and with three factors $\ddc u$ integrate to zero, because $\omega^2$ is closed and $\ddc u$ is exact. For the remaining term, put $S=\ddc\omega$. Integration by parts gives
\[
 \int_X(\ddc u)^2\wedge\omega
 =\int_Xu\,\ddc u\wedge S
 =-\int_X\mathrm{d}u\wedge\dc u\wedge S,
\]
where the last equality uses $\mathrm{d}S=0$. This proves the identity. Its last integral is nonnegative, and $u=0$ attains the upper bound $3TA$. Since smooth and bounded potentials have the same mass extrema, by \cref{thm:gaud-volume-mass-extrema}, this proves $v_+(\omega)=3TA$.

If $c_1(L)=0$ in real cohomology, then $c$ is exact and Stokes' theorem gives $\int_Bc^2=0$. The pointwise identity $c^2=-\lambda^2\eta^2$ then forces $c=0$, and $\omega$ is closed. Otherwise $S=-\pi^*(c^2)$ is nonzero, smooth and semipositive. For a pluriclosed metric $\gamma$, Stokes' theorem would give
\[
 0<\int_X\gamma\wedge S
 =\int_X\omega\wedge\ddc\gamma=0,
\]
a contradiction. Likewise, a closed positive current $R\geq\varepsilon\omega$ with $\varepsilon>0$ would give $0<\int_XR\wedge S=\int_X\omega\wedge\ddc R=0$. All these pairings are pairings of currents with smooth forms.

Finally, let $w$ be a local holomorphic fiber coordinate given by a frame of $L$, and let $\sigma$ be a local holomorphic two-form on $B$. The form $\pi^*\sigma\wedge\mathrm{d}w/w$ is invariant under the quotient action. Changing the frame of $L$ adds to $\mathrm{d}w/w$ a horizontal one-form, whose wedge product with $\pi^*\sigma$ vanishes. These local sections of $K_X$ therefore have the transition functions of $\pi^*K_B$, which proves $K_X\simeq\pi^*K_B$.
\end{proof}

\Needspace{14\baselineskip}
For separated potentials, positivity bounds both the amplitude and the downward slope of the vertical derivative. The sharp variance bound under these two constraints determines how much mass such a potential can lose. For $M,T>0$, set
\[
 F(M,T)=
 \begin{cases}
  T^2/12,&T\leq2M,\\
  M^2-4M^3/(3T),&T\geq2M.
 \end{cases}
\]

\begin{lemma}\label{lem:periodic-variance-bound}
If a smooth real $T$-periodic function $p$ has mean zero, $|p|<M$ and $p'>-1$, then
\[
 \frac1T\int_0^T p^2\,\mathrm{d}s<F(M,T).
\]
The supremum of the left side over this class is $F(M,T)$.
\end{lemma}
\begin{proof}
For $p=0$ the inequality holds because $F(M,T)>0$. Otherwise set $a=\max p>0$, $-b=\min p<0$ and $r=a+b$. Integrating $-p'<1$ over a forward interval $I$ from a maximum to the next minimum gives $r<|I|\leq T$; also $r<2M$, since $|p|<M$. The function $G(s)=(a-s)(s+b)$ is nonnegative on the range of $p$. Since $p$ has mean zero and $-p'<1$, integration over $I$ gives
\[
 abT-\int_0^T p^2\,\mathrm{d}s
 =\int_0^T G(p)\,\mathrm{d}s
 \geq-\int_I G(p)p'\,\mathrm{d}s
 =\frac{r^3}{6}.
\]
As $ab\leq r^2/4$, the average of $p^2$ is at most $r^2/4-r^3/(6T)$. This expression increases strictly for $0<r<T$, and its value at $r=\min\{2M,T\}$ is $F(M,T)$. Since $r<\min\{2M,T\}$, the strict inequality follows.

For sharpness, consider the following $T$-periodic profiles $p_0$. When $T\geq2M$, let $p_0$ be constant equal to $M$ and to $-M$ on two intervals of length $(T-2M)/2$ each, joined by a descent of slope $-1$ and length $2M$, and close the period by an upward jump. When $T\leq2M$, let $p_0(s)=T/2-s$ for $0\leq s<T$, with its upward jump at the end of the period. In both cases $p_0$ has mean zero, $|p_0|\leq M$, $p_0'\geq-1$ in the sense of distributions, and the average of $p_0^2$ is $F(M,T)$; both constructions apply when $T=2M$. Convolution with a nonnegative periodic approximate identity preserves the mean and the bound on the distributional derivative, and multiplication by $1-\varepsilon$ makes both inequalities strict. As the convolution scale and $\varepsilon$ tend to zero, the resulting smooth functions converge to $p_0$ in $L^2$, which proves that the supremum is $F(M,T)$.
\end{proof}

\Needspace{18\baselineskip}
\begin{proposition}\label{prop:primitive-elliptic-separated}\index{Monge--Ampere mass@Monge--Amp\`ere mass!separated potentials}
Suppose $c_1(L)\ne0$ in real cohomology, and put
\[
 \kappa=-\frac{\int_Bc^2}{A}>0.
\]
Every smooth potential $u=f(t)+\pi^*v$ with $f$ of period $T$, $v\in C^\infty(B,\R)$ and $\omega+\ddc u>0$ satisfies
\[
 \int_X(\omega+\ddc u)^3
 >3TA\bigl(1-\kappa F(\kappa^{-1/2},T)\bigr)>0.
\]
Let $\lambda\geq0$ be the nonnegative eigenvalue function of $c$ relative to $\eta$, so $c^2=-\lambda^2\eta^2$, and put $M_0=1/\max_B\lambda$. For the subclass $v=0$, the exact mass infimum is
\[
 3TA\bigl(1-\kappa F(M_0,T)\bigr)>0.
\]
It is approached by smooth potentials $f(t)$ with $\omega+\ddc f(t)>0$, and it is not attained by such potentials.
\end{proposition}
\begin{proof}
Put $p=f'$. The chain rule gives
\[
 \omega+\ddc u
 =\pi^*(\eta+pc+\ddc v)+(1+p')D.
\]
Positivity on the vertical line implies $p'>-1$. Over every point of $B$, the function $t$ takes every value modulo $T$, so restriction to $\ker\partial t$ shows that $\eta+pc+\ddc v$ is a K\"ahler form on $B$ for every value of $p$. Its square has integral
\[
 \int_B(\eta+pc+\ddc v)^2=A(1-\kappa p^2)>0,
\]
since primitivity removes the mixed term and Stokes' theorem removes the exact terms. Hence $|p|<\kappa^{-1/2}$.

Only the vertical derivative enters \eqref{eq:primitive-elliptic-mass}, because $-\pi^*(c^2)$ has top degree in the base directions. For a $T$-periodic function $F_0$, integration over a fiber gives $\int_E F_0(t)D=\int_0^T F_0(s)\,\mathrm{d}s$, and therefore
\[
 \int_X(\omega+\ddc u)^3
 =3TA\left(1-\frac\kappa T\int_0^T p^2\,\mathrm{d}s\right).
\]
Now apply \cref{lem:periodic-variance-bound} with $M=\kappa^{-1/2}$. The formula for $F$ gives $F(M,T)<M^2$, which proves positivity of the bound.

For $v=0$, diagonalizing $c$ with respect to $\eta$ shows that the horizontal positivity condition is exactly $|p|<M_0$; together with $p'>-1$ it also suffices for positivity of the full form. Moreover $\kappa$ is the $\eta^2$-average of $\lambda^2$, so $\kappa\leq M_0^{-2}$. \Cref{lem:periodic-variance-bound} with $M=M_0$ therefore proves the claimed positive lower bound and its nonattainment. The smooth approximating functions in its proof have mean zero and hence periodic primitives $f_\varepsilon$, and the potentials $f_\varepsilon(t)$ prove that the bound is sharp.
\end{proof}

The proof of the first bound does not show that it is sharp for separated potentials over an arbitrary base: positivity of the classes $[\eta]+p[c]$ alone does not provide one function $v$ that makes all the representatives $\eta+pc+\ddc v$ positive. The threefold of \cref{sec:vanishing-lower-mass} is the primitive elliptic bundle over $\mathbb P^1\times\mathbb P^1$ with $q=\mathrm e^{2\pi}$, so that $A=2$, $\kappa=1$ and $T=4\pi$. For it, \cref{prop:primitive-elliptic-metric,prop:primitive-elliptic-separated} reduce to \cref{prop:lower-mass-balanced-metric,prop:lower-mass-no-pluriclosed,prop:elliptic-radial-mass}, the bound is $8$ and it is sharp, and by \cref{thm:vanishing-lower-mass} its lower mass vanishes for every Hermitian form.

A base without curves rules out the averages of horizontal curves used in the known collapse example. The following construction uses the period theorem for K3 surfaces and the description of their K\"ahler cone \cite[Theorems 7.4.1 and 8.5.2 of the online version]{HuyK3}.

\begin{example}\label{ex:primitive-elliptic-k3}\index{Elliptic fibration!primitive K3 example}
There is a K3 surface $B$ with $\operatorname{Pic}(B)=\mathbb Z[L]$, $c_1(L)^2=-4$, and a K\"ahler form $\eta$ orthogonal to $c_1(L)$. It has no compact curves. The associated $X$ has trivial canonical bundle, a balanced metric of finite upper mass, and no pluriclosed metric or closed K\"ahler current. It has no divisors, and every compact irreducible curve is an elliptic fiber.
\end{example}
\begin{proof}
In a hyperbolic plane of the K3 lattice, choose a basis $e,f$ with $e^2=f^2=0$ and $e\cdot f=1$, and put $\ell=e-2f$. This vector is primitive and has square $-4$. Choose a period in $\ell^\perp$ outside the hyperplanes perpendicular to integral vectors not in $\mathbb Z\ell$. Such a period exists by Baire's theorem: the period domain in $\ell^\perp$, of signature $(3,18)$, is nonempty and each excluded hyperplane is proper. The period theorem and the description of integral $(1,1)$-classes give the asserted Picard group. An irreducible curve would have nonzero class, because its K\"ahler area is positive, and hence square $-4m^2<-2$ for some integer $m\neq0$, contradicting adjunction. Thus there are no curves.

The orthogonal complement of $\ell$ in $H^{1,1}(B,\R)$ has signature $(1,18)$. Choose a class of positive square there, with sign chosen so that it lies in the component of the positive cone containing the K\"ahler classes. Since there are no rational curves, the K\"ahler-cone theorem shows that this class is K\"ahler; choose $\eta$ in it. \Cref{prop:primitive-elliptic-metric} then gives the assertions about the metric and the canonical bundle.

If a divisor $Z\subset X$ dominated $B$, the push-forward of its integration current would be $\pi_*[Z]=m$ for a positive generic degree $m$. Stokes' theorem and the projection formula would then give
\[
 0=\int_X[Z]\wedge\ddc\omega
   =m\int_B(-c^2)>0.
\]
Thus no divisor dominates $B$. The image of any other irreducible divisor would be a curve or a point; the base has no curves, and a fiber has dimension one, so both alternatives are impossible. Likewise, the image of an irreducible compact curve is a point, and the curve is therefore an entire elliptic fiber.
\end{proof}

Put $V=3TA$. The root volumes of complementary tests can be computed explicitly for one radial family. For smooth real $T$-periodic functions $a,g$, consider
\[
 \Gamma=a(t)\eta^2+g(t)\eta\wedge D.
\]

\Needspace{15\baselineskip}
\begin{proposition}\label{prop:primitive-radial-root-volume}\index{Gauduchon volume!primitive radial tests}
The form $\Gamma$ is strictly positive and $\ddc$-closed exactly when $a$ is a positive constant and $g>0$. For such a form put $\bar g=T^{-1}\int_0^Tg(s)\,\mathrm ds$ and $d_\Gamma=\int_X\omega\wedge\Gamma$. Its positive square root and volume are
\[
 \gamma=\sqrt a\,\eta+\frac{g(t)}{2\sqrt a}D,
 \qquad
 \int_X\gamma^3=\frac32TA\sqrt a\,\bar g
       \leq\frac{d_\Gamma^{3/2}}{\sqrt V}.
\]
Equality holds exactly when $\bar g=2a$, without requiring $g$ to be constant. In particular every degree-one test in this family has root volume at most $V^{-1/2}$.
\end{proposition}
\begin{proof}
Since $\eta\wedge c=0$, the form $g\eta\wedge D$ is closed. The chain rule gives
\[
 \ddc(a\eta^2)=a''D\wedge\eta^2,
\]
the $a'c\wedge\eta^2$ term having excess base degree. Thus $\ddc\Gamma=0$ forces $a''=0$, and periodicity makes $a$ constant. In a positive coframe with $\eta=\sigma_1+\sigma_2$, $D=\sigma_3$, the complementary coefficient matrix of $\Gamma$ is $\operatorname{diag}(g,g,2a)$. Hence positivity is exactly $a,g>0$.
The displayed $\gamma$ is positive and squares to $\Gamma$. Multiplication and fiber integration give
\[
 d_\Gamma=TA(a+\bar g),\qquad
 \gamma^3=\tfrac32\sqrt a\,g\eta^2\wedge D.
\]
For fixed $d_\Gamma>0$ set $s=d_\Gamma/(TA)$. Then $0<a<s$, $\bar g=s-a$, and $\sqrt a(s-a)$ has its unique maximum at $a=s/3$. Substitution proves the bound and its equality condition.
\end{proof}

For the divisor-free threefold of \cref{ex:primitive-elliptic-k3}, \cref{thm:divisor-free-transverse-concentration} gives divergent transverse pairings against normalized Gauduchon tests. The radial family above has bounded root volume at fixed degree, even when $g$ concentrates arbitrarily strongly. Large transverse pairings do not by themselves give unbounded root volumes or a sequence of potentials with masses tending to zero.

\begin{problem}\label{prob:primitive-elliptic-lower-mass}
Determine the lower Monge--Amp\`ere mass (\cref{def:bmp}) of the primitive elliptic bundles other than the threefold of \cref{sec:vanishing-lower-mass}, for instance of their balanced metrics $\omega$ from \cref{prop:primitive-elliptic-metric}, in particular for the bundles of \cref{ex:primitive-elliptic-k3}. Does it vanish whenever $c_1(L)\neq0$ in real cohomology?
\end{problem}

The estimates for separated potentials do not answer this question; the known collapse over $\mathbb P^1\times\mathbb P^1$ uses potentials outside that class. On the K3 example every compact curve is vertical, so averaging over such curves cannot produce the horizontal directions used in \cref{lem:lower-mass-curve-average}. A proof of collapse would therefore need a different construction, while positive lower mass would require a bound for all admissible potentials or for all normalized Gauduchon tests.


\subsection{Null loci and the singularities of the Chern--Ricci flow}

Let $X$ be a connected compact complex manifold of dimension $n$, let $\theta$ be a smooth real $(1,1)$-background, and fix a smooth strictly positive probability form $\mu$. Write $\mathcal V_\mu$ for the volume in \cref{def:gaud-variational-volume}. In discussing the flow, $\omega_0$ is the initial Hermitian metric, $\operatorname{Ric}(\omega_0)$ its conventional Chern--Ricci form, and $\tau$ the finite maximal time; thus $\theta_\tau=\omega_0-\tau\operatorname{Ric}(\omega_0)$. An annihilating current $T$ means a positive $\ddc$-closed $(n-1,n-1)$-current normalized by the fixed Hermitian trace in \cref{cor:chern-ricci-current-lifetime}.

For a big smooth background $\theta$, Dang defines its non-Hermitian locus $E_{\mathrm{nH}}(\theta)$ as the intersection of the singular sets of all strict currents in that background with \emph{almost analytic singularities}; this intersection is itself attained by such a current. In the source's terminology the potential is locally a positive multiple of the logarithm of an ideal norm, up to a bounded remainder, and on a suitable log resolution its remainder after subtracting the divisorial logarithms is smooth. We retain this exact singularity condition.

\Needspace{8\baselineskip}
For nef and big $\theta$, define
\[
 \operatorname{Null}(\theta)
 =\bigcup_{\substack{V\subset X\text{ irreducible analytic}\\
                     \dim V>0,\ \theta|_V\text{ not big}}}V.
\]
Bigness on a singular subvariety uses its intrinsic analytic definition. Dang proves $E_{\mathrm{nH}}(\theta)=\operatorname{Null}(\theta)$ for arbitrary compact complex $X$ and nef, big, possibly nonclosed $\theta$. This qualitative identification is therefore a known theorem. Detecting bigness and this locus from numerical data is a further issue: the numerical criterion recorded in \cref{sec:history-geometry} uses bounded mass on $X$ and positive lower volume on its positive-dimensional submanifolds.

The known flow theorem recorded in \cref{sec:history-geometry} identifies finite-time Chern--Ricci singularities with the null locus under positive limiting lower volume, for arbitrary initial metrics in dimension two, and for Guan--Li initial metrics in higher dimensions. Here Guan--Li means $\ddc\omega_0=0$ and $\mathrm{d}\omega_0\wedge \dc\omega_0=0$. These hypotheses must not be silently replaced by positivity of $\mathcal V_\mu$.

The current calculation in \cref{cor:chern-ricci-current-lifetime} applies without those extra hypotheses. At a finite maximal time $\tau$, the background $\theta_\tau=\omega_0-\tau\operatorname{Ric}(\omega_0)$ is nef. Its normalized annihilating currents satisfy
\[
 \int_X\theta_\tau\wedge T=0.
\]
If $\mathcal V_\mu(\theta_\tau)>0$, their determinant-weighted measures vanish by \cref{prop:nef-annihilator-determinant}. This is a precise degeneracy statement, but it does not imply that their traces are concentrated on analytic sets. A smooth positive semidefinite matrix of deficient rank can be defined on an open set.

\Needspace{8\baselineskip}
\begin{problem}\label{qu:chern-ricci-gaud-singularities}
Suppose a Chern--Ricci flow has finite maximal time $\tau$ and $\mathcal V_\mu(\theta_\tau)>0$. Can one prove that $\theta_\tau$ is big and that the flow converges smoothly to a positive metric locally on $X\setminus E_{\mathrm{nH}}(\theta_\tau)$, without a bounded mass hypothesis? How do the annihilating Gauduchon currents describe the resulting degeneracy?
\end{problem}

The first assertion is a case of the volume-to-bigness question. Once bigness of $\theta_\tau$ is known, the parabolic theorem recorded in \cref{sec:history-geometry} supplies convergence as currents and smooth local convergence off $E_{\mathrm{nH}}(\theta_\tau)$, by parabolic estimates. Combined with nefness and the qualitative null-locus theorem, this identifies the analytic exceptional set. Thus the missing step in this route is the proposed numerical criterion for bigness. The local envelope continuity and homogeneous replacement arguments of \cref{thm:big-continuous-envelope-contact} supply an elliptic foundation away from a strict barrier's poles; they are not a reconstruction of those parabolic estimates.

Completeness requires an additional metric estimate. In \cref{thm:incomplete-divisorial-metric}, the unique finite-energy solution of a divisorial exponential equation in a big closed class is incomplete on every nonempty connected open subset of its ample locus. Its minimal divisorial singularity cancels the prescribed pole. This resolves the unrestricted big-class completeness question negatively; it does not decide completeness under nefness or under the compatibility conditions of the geometric smoothing problem discussed in \cref{sec:history-completeness}. In particular it does not settle \cref{qu:chern-ricci-gaud-singularities}, whose limiting background is nef.

For integrable poles weakened by logarithmic factors, \cref{cor:divisorial-exponential-completeness} gives a sufficient condition in terms of the density, the potential and the tangential determinant. The proof controls normal length and applies to arbitrary Hermitian backgrounds. The torus models in \cref{ex:critical-complete-divisor,ex:supercritical-incomplete-divisor} attain its critical exponent and show its sharpness, even with a K\"ahler background and full mass. Their pole coefficients need not extend smoothly and positively across the divisor; this distinction preserves the separate scope of the geometric smoothing question.

\subsection{Degenerate densities and uniqueness}

Let $X$ be a connected compact complex manifold, let $\omega$ be a Hermitian form and let $\mathrm dV$ be a smooth positive volume form. For $u\in\PSH(X,\omega)$ write $\mathrm{MA}_\omega(u)$ for its Monge--Amp\`ere measure. Densities $f$ below are smooth nonnegative functions on $X$; normalization of a potential means $\sup_Xu=0$. The scalar $c$ and exponential parameter $\lambda$ are positive real numbers.

The counterexample in \cref{thm:degenerate-hermitian-nonuniqueness} shows that a smooth nonnegative density can have two distinct smooth normalized solutions even on $\mathbb P^3$. The further construction in \cref{thm:numa-thin-zero-set} has uncountably many smooth normalized solutions and a density positive off two real hypersurfaces. By modifying the background, \cref{thm:quadratic-density-nonuniqueness} achieves exactly quadratic vanishing on those hypersurfaces. Thus neither zero smooth volume of the zero set nor finite-order vanishing suffices for uniqueness. All the backgrounds are strictly Hermitian and have bounded mass; their non-Hermitian loci are empty. The degeneration causing nonuniqueness belongs to the density and the solution forms, and is not detected by the background's null locus.

This places a necessary restriction on any geometric use of the equation. A construction that assigns a canonical potential to arbitrary nonnegative data must specify a selection procedure and prove that it is well defined. Existence in the ordinary minimal singularity type does not provide such a selection. The unique scalar, and the corresponding variational volume for strictly positive data, remain meaningful without uniqueness of every potential. On a compact Hermitian surface, \cref{uni-cor:dense-lowdim} proves uniqueness for $C^2$ potentials and continuous density positive on a dense set, with no differential condition on the metric. In dimension three, \cref{uni-thm:sobolev} instead uses a sign condition on $\ddc\omega$. The cofactor identity in \eqref{eq:mixed-uniqueness-defect} and the signed quadratic-difference argument provide two distinct integration mechanisms.

\begin{problem}\label{qu:density-zero-set-uniqueness}
Let $X$ be connected, $\omega$ Hermitian and $f\geq0$ smooth. Suppose $\{f>0\}$ is connected and dense. Are two bounded normalized solutions of $\mathrm{MA}_\omega(u)=cf\,\mathrm{d}V$, with the same $c>0$, necessarily equal? Separately, for densities allowing nonuniqueness, do the solutions of $\mathrm{MA}_\omega(u_\lambda)=\mathrm{e}^{\lambda u_\lambda}f\,\mathrm{d}V$ converge uniformly after subtracting their suprema, as $\lambda\downarrow0$?
\end{problem}

For $C^{1,1}$ potentials the first question is answered affirmatively by \cref{uni-thm:anchors}. The remaining issue is the passage to bounded potentials. All the hypersurface counterexamples here have three positive-density components; even the analytic density in \cref{uni-thm:analytic-example} has exactly this topology. The second asks for a canonical limit, and does not follow from uniqueness of the positive-exponent equation at each fixed parameter. A selection could exist even when the limiting equation has several solutions. These counterexamples do not refute \cref{conj:continuous-gaud-adjoint}: that conjecture uses a finite continuous logarithmic density $F$, for which $\mathrm{e}^F$ is strictly positive. The logarithms of the degenerate densities have no finite continuous extensions across their zero sets.

There is already some control of selection in the Hermitian $L^p$ setting: \cref{prop:symmetric-equation-clusters} proves uniform precompactness of the normalized approximations, and every cluster limit preserves the symmetries of the data. For the symmetric projective example, the radial argument of \cref{thm:numa-selected-limit} goes further: the whole normalized family converges uniformly to the explicit symmetric $C^{1,1}$ solution of \cref{prop:numa-radial-continuum}. Thus this example has both many normalized solutions and a definite limit for the specified approximation. The general selection question above remains open in this treatment.

For a nonzero density at positive threshold in dimension at least two, \cref{cor:smooth-dual-exact-equation,cor:smooth-dual-exponential-limit} give uniqueness of the normalized threshold solution and uniform convergence of the whole normalized positive-exponent family whenever a smooth strictly positive maximizing Gauduchon current exists. The construction in \cref{ex:nonclosed-dual-zero-density} satisfies this condition even though the density vanishes on an open set. Thus the size of the zero set alone does not determine uniqueness; the dual condition also involves the background.

\subsection{Balanced and pluriclosed metrics}

Let $X$ be a connected compact complex manifold of dimension $n\geq2$. The forms $\beta$ and $\omega$ below are Hermitian metrics on $X$, and $c_1(X)$ denotes the first real Chern class of its holomorphic tangent bundle. The Fujiki class $\mathcal C$ consists of manifolds bimeromorphic to compact K\"ahler manifolds.

A Hermitian metric $\beta$ is balanced if $\mathrm{d}\beta^{n-1}=0$, while a Hermitian metric $\omega$ is pluriclosed if $\ddc\omega=0$. The Fino--Vezzoni conjecture asks whether a compact complex manifold admitting both types of metrics is K\"ahler; the two metrics can be different. Its source is recorded in \cref{sec:history-geometry}. A single metric satisfying both conditions is already K\"ahler, which does not resolve the question with two metrics.

Chiose proves that a compact manifold in $\mathcal C$ admitting a pluriclosed metric is K\"ahler. Consequently, in the balanced--pluriclosed setting, it would suffice to construct a closed big $(1,1)$-current and apply the characterization of $\mathcal C$. This route emphasizes a geometric existence problem which is absent when one only works with an arbitrary Hermitian background.

There is an elementary obstruction to using the scalar equation alone. For every distributional potential $u$,
\[
 \mathrm{d}(\omega+\ddc u)=\mathrm{d}\omega.
\]
Thus if the starting pluriclosed metric is nonclosed, no potential in its fixed background makes it K\"ahler. The finite-volume equation, including its smooth solution, does not change this differential datum. Similarly, a conformal Gauduchon representative need not be balanced. Positivity and a scalar normalization cannot be used to replace either missing closedness condition.

The balanced form $\beta^{n-1}$ nevertheless supplies more than a Gauduchon test: it pairs trivially with every exact real two-form by Stokes. It therefore interacts with flows whose variation has an exact component. Fino--Vezzoni propose long-time existence and convergence to a K\"ahler metric for the pluriclosed flow when a balanced metric exists and $c_1(X)=0$. Their invariant examples illustrate the use of a conserved balanced trace. In general that integral is not a pointwise metric bound.

\begin{problem}\label{qu:balanced-pluriclosed-program}
Use the coexistence of a balanced metric and a pluriclosed metric to construct closed positive data, or to control a geometric flow that can change the torsion. Determine what estimate beyond the finite balanced trace prevents metric degeneration. Can determinant estimates for Gauduchon currents contribute to that estimate?
\end{problem}

This asks for an interaction with the scalar theory, not an implication already proved by its equation. One can first study families for which the balanced trace and the determinant density are explicitly controlled. The passage from an integral trace bound to local ellipticity is then a concrete analytic target. No such passage for arbitrary families is asserted here.

\subsection{Morse inequalities and the choice of volume}

For the closed comparison, let $X$ be a connected compact K\"ahler $n$-fold and let $\alpha,\beta$ be closed nef real $(1,1)$-classes. Write $\operatorname{Vol}$ for the usual volume of a closed class. In the Hermitian discussion, $X$ is only assumed connected, compact and complex; $\theta,\eta$ are smooth real $(1,1)$-forms, and $\mu$ is a fixed smooth strictly positive probability form. Use $\lambda_\mu$, $\mathcal V_\mu$ and $D_\mu$ from \cref{def:gaud-variational-volume,sec:gauduchon-collapse}.

The closed K\"ahler theory provides a demanding comparison. Tosatti's recent preprint establishes the transcendental Morse inequality
\[
 \operatorname{Vol}(\alpha-\beta)
 \geq\int_X\alpha^n-n\int_X\alpha^{n-1}\wedge\beta
\]
for closed nef classes on compact K\"ahler manifolds, together with the resulting volume differentiability and cone-duality consequences. We record it as a recent literature result; its proof is not reconstructed here. The unrestricted non-K\"ahler extension remains a different problem. The qualitative bounded-mass criterion and the source of the closed theorem are recorded in \cref{sec:history-geometry}.

For nonclosed backgrounds, the displayed intersection expression is not generally invariant under changing representatives. Also, the upper class volume and $\mathcal V_\mu$ have different definitions and different possible finiteness behavior. A Hermitian analogue must first specify which invariant it estimates.

\Needspace{10\baselineskip}
The first-variation formula gives a finite quantity to investigate. Fix a real number $t>0$. Suppose $\theta_s=\theta-s\eta$ has a smooth positive representative for $0\leq s\leq t$, and let $\gamma_s$ be a minimizing Gauduchon metric for its smooth equation. By \cref{prop:gaud-volume-first-variation} and the fundamental theorem of calculus,
\begin{equation}\label{eq:gaud-volume-subtraction-cost}
 \lambda_\mu(\theta-t\eta)
 =\lambda_\mu(\theta)
   -\int_0^t
       \frac{\int_X\eta\wedge\gamma_s^{n-1}}
            {D_\mu(\gamma_s)}\,\mathrm{d}s.
\end{equation}
The quotient is unchanged by rescaling the minimizer. This identity is valid on the stated positive segment and uses no upper Monge--Amp\`ere mass bound. A uniform upper bound on its integrand gives a lower bound for the remaining root volume. Proving a useful bound near a degenerate endpoint is an additional problem; concavity alone provides the tangent upper bound, not the lower bound needed here.

\begin{problem}\label{qu:hermitian-morse-program}
Find geometric hypotheses controlling the cost in \eqref{eq:gaud-volume-subtraction-cost} up to a nef endpoint, and determine whether the resulting positive root volume forces bigness there. For closed backgrounds, compare the resulting estimate with classical Morse inequalities. For nonclosed backgrounds, retain the differential data and state explicitly any dependence on $\mu$.
\end{problem}

The singularity and divisorial constructions in the book become relevant at this stage: they record how much positivity can be spent on imposed poles, while the residual Gauduchon volume stays finite. They do not yet give the moving intersection product or the birational volume formula of the closed theory. Recovering such geometric information is a further objective, with the finite-divisor and Hopf calculations serving as explicit tests.


\section{Closed classes and LCK volume equations}

Let $X$ be a connected compact complex manifold. The closed classes in this section belong to $H^{1,1}_{\mathrm{BC}}(X,\mathbb R)$. For an LCK structure, let $\lambda$ denote its closed real Lee form and let $D_\lambda$ be the twisted second-order operator of \eqref{eq:lck-gauge-operators}. The map $j_X:H^1_{\mathrm{dR}}(X,\mathbb R)\to H^{1,1}_{\mathrm{BC}}(X,\mathbb R)$ sends $[\lambda]$ to $\{\dc\lambda\}$, as in \cref{sec:lee-gauduchon-cone}.

The loci of \cref{chap:bott-chern-loci} distinguish small negative errors for closed currents from strict closed positivity. Any proposed numerical restricted-volume construction in the general Hermitian setting must first pass \cref{prop:bc-zero-class-infinite-error-volume}: an unrestricted supremum over smooth Hermitian-error representatives is already infinite in the zero class on the Hopf threefold. The failure occurs at each fixed error size. Taking the error to zero afterwards cannot recover a class invariant from that expression.

The global subsolution threshold also constrains divisorial singularities. For a subsolution at this coefficient, \cref{prop:threshold-divisorial-constraint} places every smooth maximizing Gauduchon sequence in the annihilator of each component class in the limit. If one such sequence has uniformly integrable determinant-weighted densities, \cref{cor:semipositive-divisor-exclusion} excludes classes with suitable smooth semipositive representatives, including the hyperplane divisor on the Hopf threefold. For a nef component class, \cref{cor:nef-divisor-threshold-volume} forces its fixed-density variational volume to vanish. Root tails suffice for both conclusions. Characterizing the remaining possible components requires a realization theorem. The smooth-test closure can be smaller than the full current set, as \cref{ex:gauduchon-current-smoothing-obstruction} shows.

A Hermitian numerical characterization of the loci in \cref{chap:bott-chern-loci} requires both a finite restricted invariant related to traces and a way to turn its positivity into a strictly positive current. Neither requirement follows from local regularization alone.

The LCK setting gives a separate set of equation questions. By \cref{cor:lck-no-closed-big-class}, a compact strictly LCK manifold has no big ordinary Bott--Chern class. Its positive LCK form belongs instead to a twisted class, whose representatives are described using $D_\lambda$. The coupled system in \cref{prop:lck-coupled-equation} identifies the additional conformal equation imposed by Gauduchon's condition. The uniqueness problem \cref{q:lck-gauduchon-volume} concerns that system with a fixed Lee class and twisted class; the ordinary-potential nonuniqueness examples do not decide it.

Two geometric questions use the Lee form in different ways. The pseudoeffectivity conjecture \cref{conj:lck-lee-psef} requires positivity against all Gauduchon tests, extending the positive pairing obtained from the LCK metric itself. The Calabi--Yau conjecture \cref{conj:lck-calabi-yau-vaisman} asks for a Vaisman metric on the manifold. The proved rigidity theorem \cref{thm:lck-chern-ricci-rigidity} applies when one metric already satisfies both Chern--Ricci flatness and the Gauduchon condition. Arranging those conditions simultaneously is the remaining geometric existence issue.

The Lee--Gauduchon cone makes one part of this geometry invariant under bimeromorphic changes. Its dual is the cone of degree-one classes whose image under $j_X$ is pseudoeffective; the strict inequalities exclude $\ker j_X$. The full duality and modification argument are proved in \cref{sec:lee-gauduchon-cone}. A nonzero positive class in this image obstructs balanced metrics throughout the bimeromorphic class. The Hopf calculation gives a positive half-ray, despite unbounded Monge--Amp\`ere masses. Thus this cohomological invariant addresses a geometric obstruction while leaving the nonlinear estimates for the scalar equation as a separate task. The converse to its balanced-metric consequence fails: \cref{ex:lee-calabi-eckmann} has zero Lee--Gauduchon cone and no strongly Gauduchon metric.

\section{Degeneration orders and finite-distance boundaries}

Let $X$ be a connected compact complex $n$-fold and let $\theta$ be a nef smooth real $(1,1)$-background. Fix a Hermitian form $h$ and a smooth positive probability form $\mu$. The quantities $\underline a(\theta)$ and $\overline a(\theta)$ are the lower and upper limits of $\log\mathcal V_\mu(\theta+\epsilon h)/\log\epsilon$ as $\epsilon\downarrow0$. The metric on positive background classes is the Hessian metric of \cref{sec:gauduchon-metric-geometry}.

For a nef background the two orders $\underline a(\theta)$ and $\overline a(\theta)$ of \cref{prop:nef-degeneration-orders} are already independent of the smooth probability measure and the Hermitian perturbation direction. Their equality is not proved. Neither is integrality of a common value, if it exists. These are the precise questions in \cref{qu:nef-degeneration-exponent}.

The explicit examples give several tests. Closed nef classes on a K\"ahler manifold have the polynomial order given by numerical dimension. On the Hopf threefold, the horizontal form has order one, although every positive perturbation has infinite upper mass. On the elliptic threefold, \cref{prop:elliptic-boundary-volume} gives orders one, two and three on the nonnegative closed-class quadrant; for integral classes they equal three minus the section-growth degree. The finite-divisor example gives an explicitly quadratic decay as a coefficient reaches its feasibility threshold. A proposed general exponent must reproduce all these cases without first imposing bounded mass.

The background-class metric yields a different kind of boundary information. On a finite-dimensional cone it splits as a product of the logarithmic volume direction and a unit-volume hypersurface. Consequently a path whose volume tends to zero or infinity has infinite length. Yet the blowup example in \cref{sec:gauduchon-metric-geometry} has a finite-distance boundary with positive nef volume. It follows that completing this metric cone cannot simply mean adjoining all nef classes, or declaring every nef boundary infinitely far away.

\begin{problem}\label{qu:gaud-metric-boundary}
For a specified finite-dimensional cone of smooth background classes, describe the metric completion of its unit-volume hypersurface. Which boundary points admit a current representative, and what estimates on the minimizing Gauduchon metrics distinguish finite from infinite distance?
\end{problem}

The cone must be specified in this question. Allowing additional directions can shorten path distances. The finite-dimensional metric calculations do not establish a distance on the entire infinite-dimensional quotient.

\section{Weak Gauduchon adjoints}

Use a connected compact complex $n$-fold $X$, a smooth positive volume form $\mathrm dV$, a smooth semipositive big background $\theta$ and a model reference $\varphi$ with the smaller-background property of \eqref{eq:normalization-reference}. For $F\in C(X,\mathbb R)$, $B_\varphi(F)$ is the logarithm of the scalar in $\mathrm{MA}_\theta(u)=\mathrm e^{B_\varphi(F)+F}\mathrm dV$, with $u\simeq\varphi$. The supporting probabilities form $\mathcal A_\varphi(F)$, as defined in \cref{chap:normalization}. The conjecture below specializes to a Hermitian form $\omega$ and $\varphi=0$.

The scalar normalization provides a route to weak adjoint measures that does not require multiplying an unbounded Monge--Amp\`ere product by an arbitrary density. For the logarithmic normalization $B_\varphi(F)$, the set $\mathcal A_\varphi(F)$ of supporting probabilities is nonempty, convex, weakly compact and absolutely continuous. Smooth data in a Hermitian background give the classical Gauduchon adjoint probability. The open issue in the present treatment is whether this probability remains unique for weak data.

\begin{conjecture}\label{conj:continuous-gaud-adjoint}
For a Hermitian background $\omega$, the nonsingular reference $\varphi=0$, and every real $F\in C(X)$, the set $\mathcal A_0(F)$ consists of a single probability measure.
\end{conjecture}

Equivalently, the concave normalization functional is G\^ateaux differentiable at every continuous logarithmic density in this regime. The conjecture holds for smooth $F$ by the linearized equation, and in the closed K\"ahler setting by the explicit normalization. It would imply that smooth approximations of $F$ have a unique weak limit of their Gauduchon adjoint probabilities, by \cref{prop:weak-adjoint-stability}.

A direct bounded-potential result is \cref{prop:bounded-adjoint-persistence}. If every smooth positive potential perturbation of $\omega$ is Gauduchon, normalized Monge--Amp\`ere measures extend the smooth adjoint probabilities and converge under capacity convergence of uniformly bounded potentials. This applies on every Gauduchon surface. It does not define a weak linearized operator or settle approximation by unbounded potentials in a general background.

There is also a partial result from Baire category: \cref{thm:generic-adjoint-uniqueness} proves uniqueness on a dense $G_\delta$ of continuous logarithmic densities, in every fixed background and reference regime of \cref{chap:normalization}. At those densities all choices of approximating supporting probabilities have the same weak limit. The assertion is generic in the uniform topology; it does not eliminate exceptional continuous densities.

The compact current formulation gives a more precise target. By \cref{thm:gaud-current-all-adjoints}, the supporting probabilities are exactly the normalized determinant-weighted measures of maximizing Gauduchon currents. Thus the conjecture asks whether these weighted measures agree for all maximizers. \Cref{cor:gaud-adjoint-lq} already bounds their densities uniformly in every $L^q$, $1<q<n/(n-1)$, on bounded sets of $F$; the endpoint is available if $v_-(\omega)>0$. These estimates give compactness, but do not prove uniqueness. The failure of the endpoint for all normalized Gauduchon currents on the lower-mass example does not imply failure for the maximizing currents associated with one fixed density. Another approach would show directly that the upper and lower directional derivatives of $B_0$ agree.

For a singular reference or a degenerate background, the same uniqueness question remains meaningful with the hypotheses of \cref{chap:normalization}. A further geometric interpretation requires a well-defined weak linearization. For a potential $u\in\PSH(X,\theta)$ write $\theta_u=\theta+\ddc u$, and let $\rho$ be a measurable weight. The expression $\ddc(\rho\theta_u^{n-1})$ is not a definition when multiplication of the product by $\rho$ has not been justified.

\section{A topology that retains singularities}

Let $X$ be a connected compact complex manifold, let $\omega$ be Hermitian, and let $\theta$ be a smooth real $(1,1)$-background. For $u\in\PSH(X,\theta)$, $P_\theta[u]$ and $P_\theta[u]_{\mathcal I}$ are its ordinary and multiplier-ideal envelopes from \cref{def:singularity-envelope,def:hermitian-i-envelope}. Use smooth positive volume for $L^1(X)$.

The complete metric supplied by Gauduchon traces describes normalized potentials with their $L^1$ topology. The ideal-envelope metric describes multiplier-ideal classes through their canonical representatives. Both are available without bounded mass. Neither gives the desired analogue of the source singularity geometry. The dyadic-envelope criterion of \cref{cor:dyadic-goodness-criterion} now characterizes $P_\theta[u]=P_\theta[u]_{\mathcal I}$ along a specified analytic regularization, without BMP. It does not define convergence for arbitrary ordinary singularity types; analytic classes are dense in the entire elementary ideal metric by \cref{cor:fixed-hermitian-analytic-density}.

The curve completion theorem is a sharp diagnostic. Ideal envelopes with finite logarithmic poles and zero non-pluripolar mass are dense in every normalized potential, including smooth potentials. Thus completing the ideal-envelope metric can add many potentials with identical multiplier ideals. At the same time, convergence of these trace measures does not imply capacity convergence of the potentials. A construction based solely on weak trace convergence will therefore lose nonlinear information needed by the equation.

\Cref{prop:positive-mass-ideal-discontinuity} shows that even a fixed positive mass can persist along ideal envelopes converging in the elementary ideal metric to a smooth potential of larger mass. A lower bound on the individual masses therefore does not restore mass continuity for that metric.

On a curve, \cref{thm:curve-atomic-source-metric,cor:curve-ideal-source-completeness} identify the source topology on positive-mass ideal envelopes with total variation of their atomic data. A uniform positive mass bound gives completeness in that topology, while the space is not totally bounded. This supplies a concrete comparison for the general problem.

\begin{problem}\label{qu:singularity-topology-program}
Construct a specified convergence for singularity types on an arbitrary Hermitian manifold that characterizes $P_\omega[u]=P_\omega[u]_{\mathcal I}$ by approximation with analytic singularities. State the normalization, the allowed curvature loss of the approximants, and the behavior under smooth changes of background. Determine whether this convergence comes from a metric, and prove any claimed triangle inequality or completeness.
\end{problem}

Finite damped measures provide analytic information throughout the potential space, including infinite-mass potentials. They are insensitive to pluripolar mass and their continuity follows already from capacity convergence. Thus a criterion formed only by adding damped weak convergence to capacity convergence does not add control of multiplier ideals. A plausible construction should also specify ideal inclusions at each fixed weight, or equivalent divisorial data on all modifications. Whether these conditions have the required sufficiency is a problem, not a theorem asserted here.

The positive-volume characterization in \cref{prop:conditional-i-characterization} is a useful test in the bounded-mass special regime. A general construction should recover that result where its hypotheses hold, while still making sense for the infinite-mass envelope-good potential of \cref{prop:goodness-infinite-mass}.

\Needspace{17\baselineskip}
\section{Exact trace potentials in higher dimension}

Let $X$ be a connected compact complex manifold, let $Y\subset X$ be a connected compact complex submanifold of positive dimension, and fix a quasi-plurisubharmonic function $u$ with generic Lelong number $\nu(u,Y)=0$. For a prime divisor $E$ on a smooth compact modification $\pi:Z\to Y$, $\tau_{Y,u}(E)$ is the restricted-ideal trace number of \eqref{eq:trace-profile}. For a sufficiently large Hermitian form $\Omega$ on $Y$, $\mathcal T_\Omega(u;Y)$ is the envelope in \eqref{eq:trace-profile-envelope}; $\nu(a,E)$ denotes the generic Lelong number of a potential $a$ along $E$.

The restricted multiplier ideals determine an intrinsic additive divisorial profile. Analytic ambient singularities have exact ordinary restrictions, and \cref{thm:curve-fibration-exact-trace} now gives exact higher-dimensional traces over an ambient curve problem. The latter have an ordinary type independent of a sufficiently large background. For arbitrary data, \cref{thm:trace-profile-model-envelope} still gives a model envelope of positive finite mass whose orders dominate the profile.

\Needspace{10\baselineskip}
\begin{problem}\label{qu:exact-trace-program}
In the setting of \cref{thm:trace-profile-model-envelope}, can the background be chosen so that
\[
 \nu\bigl(\mathcal T_\Omega(u;Y)\circ\pi,E\bigr)
       =\tau_{Y,u}(E)
\]
for every prime divisor $E$ on every smooth modification $\pi$ of $Y$? If so, determine a background-compatible ordinary or ideal singularity class of this realization, and prove its addition and approximation properties.
\end{problem}

The quantified equality is stronger than realizing the divisorial part on one model. \Cref{cor:divisorial-current-ordinary-type} now gives a canonical ordinary type for that part by combining its Gauduchon mass bound with a smooth current decomposition. It does not determine higher-codimensional trace data on further modifications. The exceptional-divisor formula in \cref{thm:trace-exceptional-difference} now describes the discrepancy between the canonical types on two models, including a cocycle and approximation in linear mass. \Cref{ex:exceptional-trace-with-exact-potential} shows that a nonzero exceptional term is compatible with an exact trace potential: the term records information absent from the original model's divisorial part. Increasing the background ensures a competitor with controlled curvature; it does not automatically exclude additional divisorial orders. The curve-fibration proof works because every further-model order reduces to an explicit one-variable multiplicity. A general proof needs another source of exactness.

Even exact realization does not control the mass of all ideally equivalent representatives. In \cref{prop:exact-trace-infinite-representatives}, a product with the Hopf threefold has one smooth positive extended measure and one infinite-mass representative with exactly the same ideals and trace numbers. A proposed volume or continuity characterization must specify which representative or mass domain it uses.

\section{Compatible divisors and recovery from perturbations}

Finite divisors on a fixed manifold now have an exact Gauduchon feasibility threshold and a residual variational volume. If the residual background has a smooth positive representative and the prescribed probability measure is smooth and strictly positive, \cref{thm:divisorial-residual-equation} gives a smooth residual solution. Every threshold subsolution whose divisorial coefficients dominate the prescribed ones differs from it by a constant; no bounded-remainder assumption is needed. A solution in the prescribed ordinary type therefore has the same scalar and a smooth remainder. Divisorial currents on each modification also have finite Gauduchon mass. These facts are compatible with the lack of any uniform upper Monge--Amp\`ere mass bound on the original manifold.

Passing to all modifications introduces two further requirements. A smooth positive probability measure can pull back to a density that vanishes on the exceptional divisor. The smooth dual formula cannot be applied to that density without proving a degenerate extension. The residual classes can now be compared explicitly: \cref{cor:affine-trace-exceptional-difference} subtracts the effective exceptional Bott--Chern class while retaining the prescribed exterior derivative. The coefficient approximants converge in linear trace mass, but no nonlinear volume transformation follows from this formula. Neither the degenerate equation issue nor the volume issue is resolved by assigning an intersection number to a nonclosed form.

The envelope-system construction in \cref{chap:hermitian-envelope-systems} gives an order adjunction and an exact criterion for recovery: the perturbation limit must retain the prescribed ideal type. Its ruled-surface example shows that additional divisorial singularities can be forced in a limit with zero residual volume. This is an obstruction to unrestricted recovery, not a counterexample to a theorem whose hypotheses require positive volume.

\Needspace{9\baselineskip}
\begin{problem}\label{qu:positive-system-recovery}
Give a sufficient positive-volume condition on a compatible Hermitian envelope system that forces its perturbation limit to preserve the ideal type at each parameter. Specify which volume is used, its transformation under modifications, and the uniformity required in the perturbation parameter.
\end{problem}

For fixed finite divisor support, \cref{prop:finite-divisorial-equation-limits} gives a positive-cone case: if the coefficients converge and the limiting residual background still has a smooth positive representative, the normalized smooth remainders converge smoothly. The singular equation solutions converge in capacity, their measures converge in total variation, and their logarithmic coefficients remain exact on every modification. Multiplier ideals can nevertheless jump at an integer coefficient. This concerns the canonical equation family on one manifold; recovery of a general envelope system and estimates uniform over infinitely many models still require additional arguments.

\section{Sections, kernels and the dimension of growth}

A Hermitian form on the manifold does not become the curvature of a holomorphic line bundle unless it is closed. The algebraic part of the theory must therefore keep its closed curvature data and its analytic background separate. The weighted spaces of sections, their multiplicative variant, and their intrinsic flag bodies do so. Their growth dimension can be strictly less than the complex dimension of the manifold.

For partial kernels, the general upper envelope bound is proved. A matching lower bound remains a separate problem. \Cref{thm:fujiki-partial-kernel-limit} now gives the full-dimensional kernel measure limit on a manifold with a K\"ahler modification and a positive-mass closed-curvature potential. The proof retains the relative canonical bundle: its Jacobian factor corrects both ideal integrability and the Hilbert norm. This geometric assumption implies BMP; the birational transformation identities themselves do not require it. On a holomorphic submersion with descended data, the source equilibrium theorem on the base gives a different exact limit upstairs. Write $n$ for the complex dimension of the Hopf manifold and $m$ for the positive integer tensor power of the selecting line bundle. The Hopf example has growth degree $n-1$, and normalizing by $m^n$ would erase it. Adding the vertical Hermitian form recovers a nonzero measure, with the factor computed in \eqref{eq:hopf-partial-equilibrium}.

Isolated point poles which do not descend are now treated in \cref{thm:isolated-pole-kernel-limit}: their projected jet conditions differ from the base multiplier ideals by only finitely many orders. For a single point on the Hopf manifold, \cref{thm:hopf-point-potential-limit} proves the potential lower asymptotic as well as total-variation convergence of the kernel measures. Moving the pole along its fiber leaves every selected space unchanged. These limits therefore recover the projected equilibrium, while the original Hermitian multiplier-ideal type can be different.

\Needspace{8\baselineskip}
\begin{problem}\label{qu:partial-kernel-lower-program}
Identify hypotheses on a weighted section series and its intrinsic growth dimension under which the partial Fubini--Study potentials attain the envelope lower bound and the normalized kernel measures converge. Relate the limiting measure to the Hermitian background while retaining the closed curvature and the growth dimension of the series.
\end{problem}

Beyond the isolated-point and descended regimes proved above, this question includes two tasks: producing enough sections with the required ideal constraints, and controlling the resulting measure. A Bergman upper bound supplies neither. In a non-Moishezon setting, ambient-degree normalization may already force every kernel measure to vanish, independently of singularity theory.

Regularity of extremal functions is also an active and partly solved subject. For Hermitian backgrounds, \cite[Theorems 1.1--1.2]{Ahn26} establishes locality of H\"older regularity and regularity of weighted extremals under a uniform density-in-capacity condition on the compact set. These results should be used in their stated regimes. They do not assert regularity of every singular envelope or remove the lower-asymptotic problem above.

The finite-codimension comparison suggests a more specific question on a fibration. It asks for an analytic description of the section constraints after descent.

\Needspace{18\baselineskip}
\begin{problem}\label{qu:fiber-selected-approximation}
Let $\pi:X\to B$ be a surjective holomorphic submersion of connected compact complex manifolds with connected fibers, where $B$ is K\"ahler of dimension $q\geq1$. Fix a Hermitian form $\omega$ on $X$, a holomorphic line bundle $A$ with a smooth Hermitian metric $h$ on $B$, and $\varphi\in\PSH(X,\omega)$. For positive integers $m$, write $\mathcal I(m\varphi)$ for the multiplier ideal of $\mathrm e^{-m\varphi}$. Set $\theta=c_1(A,h)$ and identify
\[
 V_m=H^0\bigl(X,\pi^*A^m\otimes\mathcal I(m\varphi)\bigr)
          \subseteq H^0(B,A^m)
\]
by holomorphic descent. Find geometric or analytic conditions that produce $\psi\in\PSH(B,\theta)$ of positive non-pluripolar mass such that, for all sufficiently large $m$,
\[
 \begin{gathered}
 W_m=H^0\bigl(B,A^m\otimes\mathcal I(m\psi)\bigr)\subseteq V_m,\\
 \dim V_m-\dim W_m=o(m^q).
 \end{gathered}
\]
Determine which features of the singularities of $\varphi$ along the fibers suffice for this approximation.
\end{problem}

For descended continuous weights and probability norms, \cref{lem:kernel-codimension-comparison} and the base equilibrium theorem would then give the kernel-measure and determinant-functional limits. They do not by themselves give convergence of the potentials. The finite isolated-point case in \cref{thm:isolated-pole-kernel-limit} supplies the stronger error $O(m^{q-1})$, under its stated hypothesis on the projected poles. \Cref{thm:horizontal-center-descent} gives the same error for a logarithmic singularity along a section over a smooth base divisor, under the curvature margin stated there. Its base reference is explicitly divisorial, and \cref{ex:horizontal-center-sharp-error} proves sharpness of the error order. The general problem does not assert that every Hermitian-positive selecting weight admits such a base reference.

\subsection{Growing vanishing orders on curves}

Fix an integer $r\geq1$ and nonnegative integers $q_m$, indexed by positive integers $m$. The model in \cref{sec:degenerate-selected-kernels} treats the smooth weight $|z|^{2r}$ and selecting ideals $(z^{q_m})$ on $\mathbb C$. When $q_m\to\infty$, the kernel vanishes at the natural rescaled limit, but has a Gaussian transition at the larger radius $(q_m/(rm))^{1/(2r)}$. For $q_m=\lfloor m\lambda\rfloor$ with a real $\lambda>0$, the limiting kernel measure omits the pole atom of the limiting envelope current. For a compact curve with the curvature and local expansion in \cref{thm:compact-growing-jets}, that theorem transfers the Gaussian transition to global sections when $q_m=o(m^{1/(2r+1)})$. The next problem concerns larger growing orders.

\begin{problem}\label{prob:growing-jets-localization}
Let $C$ be a connected compact Riemann surface, let $x\in C$, and let $(L,h)$ be a smooth Hermitian holomorphic line bundle whose curvature is semipositive and strictly positive on $C\setminus\{x\}$. Fix an integer $r\geq1$. Suppose that in a local coordinate $z$ centered at $x$ and a holomorphic frame $e$ one has
\[
 |e|_h^2=\mathrm{e}^{-\phi},\qquad
 \phi(z)=|z|^{2r}+O(|z|^{2r+1})
\]
Let $\nu$ be a smooth positive volume form with local expression $a(z)\,\mathrm{d}A$ and $a(0)=1$, where $\mathrm{d}A$ is Euclidean area. For positive integers $m$, choose integers $q_m\to\infty$ with $q_m=o(m)$. Give
\[
 V_m=H^0(C,L^m\otimes\mathfrak m_x^{q_m})
\]
the norm $\|s\|^2=\int_C|s|_{h^m}^2\,\mathrm{d}\nu$, where $\mathfrak m_x$ is the ideal sheaf of $x$. Write $B_m=\sum_j|s_{m,j}|_{h^m}^2$ for an orthonormal basis $(s_{m,j})$ of $V_m$, with value zero when $V_m=0$. Determine the growth ranges of $q_m$ for which, for each fixed real $\tau$, every sequence of points $x_m$ in the coordinate neighborhood satisfying
\[
 m|z(x_m)|^{2r}=q_m/r+\tau\sqrt{q_m/r}
\]
obeys
\[
 \frac{\pi B_m(x_m)}{mr^2|z(x_m)|^{2r-2}}
 \longrightarrow
 \Phi(\tau):=\frac1{\sqrt{2\pi}}\int_{-\infty}^{\tau}
                         \mathrm{e}^{-s^2/2}\,\mathrm{d}s.
\]
When this transition law fails, determine how further terms of the local weight change its center and scale.
\end{problem}

For the data of this problem, put $R_m=(q_m/(rm))^{1/(2r)}$. The range $q_m=o(m^{1/(2r+1)})$ is settled by \cref{thm:compact-growing-jets}: on every fixed multiple of the disc of radius $R_m$, the local weight error satisfies $mR_m^{2r+1}=q_mR_m/r=o(1)$. Its proof constructs global sections with the growing vanishing condition. Extending the range requires control of the weight remainder at the transition scale, rather than only a fixed-order kernel limit. Higher-dimensional analogues and transfer through Hermitian fibrations remain separate questions.

\section{Equations on singular spaces}

\Cref{app:hermitian-spaces} extends potentials, envelopes, ideal equivalence through resolutions, and extended and damped measures to reduced irreducible compact unibranch spaces. Its descent theorem retains the original nonclosed background. It does not assume that the resolution is K\"ahler.

The foundation does not by itself give an equation theorem on every such space. A pulled-back Hermitian form is degenerate, and a pulled-back measure may have a singular or vanishing density. Recent regularity results, including \cite[Theorem B]{Dang26}, address continuity for bounded solutions on resolutions under additional geometric and integrability hypotheses, notably the log terminal setting. Such a statement is not a result for arbitrary unibranch singularities.

The existence of singular Gauduchon metrics is already known in greater generality. For a compact reduced irreducible Hermitian space of dimension at least two, \cite[Theorems A and A']{Pan25} constructs on the regular locus a Gauduchon metric conformal and uniformly comparable to the given metric, unique up to scale within this class. The construction controls the conformal factors of suitable metrics on a resolution. This known theorem is relevant to a singular trace theory; it is not reconstructed in the present appendix and does not itself supply the nonlinear equation results asked for below.

\begin{problem}\label{qu:singular-space-equations-program}
Determine geometric and measure conditions on a compact unibranch Hermitian space that give existence, uniqueness and descent of solutions in specified singularity classes. Identify the conditions guaranteeing finite mass and continuity, and make the resulting equation independent of the chosen resolution.
\end{problem}

A useful first step is to treat a fixed resolution with explicit control of the pulled-back density and a strict big barrier. Only after that estimate is established should it be compared under a common refinement. The linear trace construction also needs a separate singular-space formulation; there is no smooth elliptic Green operator across the singular locus without additional analysis.

\section{Order of the next arguments}

Further equation work should address positive individual thresholds without positive lower mass, and descent through singular spaces. The proved general-measure results retain their two distinct geometric regimes. For bounded solutions, \cref{thm:local-big-bounded-continuity} gives local continuity wherever a fixed strict barrier is locally bounded, without a bounded mass hypothesis.

The next structural tasks are to control degeneration of normalized Gauduchon tests in the positivity conjecture, extra divisorial orders in higher-dimensional trace envelopes, and ideal-type preservation for compatible residual envelope systems with positive volume.
